\subsection{The line-tangency calculation}
\label{subsec:geometric-polynomial}

We now compute the invariant obstruction to a degree-six
arithmetic derivation preserving one nonzero homology line.
All vector spaces and maps in this subsection are over a
field of characteristic zero. The formulas are rational,
so they apply over every $\mathbb Q_\ell$.
Write $V=V_g$, $\mathbb L(V)$ for its free Lie algebra, and
$e_{\operatorname{Sp}}$ for the canonical projection onto
the symplectic-invariant part of a finite-dimensional
rational symplectic representation.

\begin{lemma}\label{lem:geometric-line-tangency}
Let $g\geq10$, let $0\ne a\in V$, and let
$D\in\mathfrak h_{g,1}(6)$ satisfy
\[
 D(a)\in[\mathbb L_6(V),a],\qquad [D,J_{g,1}]\subseteq J_{g,7}.
\]
Then $e_{\operatorname{Sp}}D\in J_{g,6}$.
Consequently every symplectically trivial quotient of a
symplectic submodule containing $D$ and annihilating its
intersection with $J_{g,6}$ sends $D$ to zero.
\end{lemma}

\begin{proof}
Choose $b$ with $\langle a,b\rangle=1$ and put
$W=\langle a,b\rangle^\perp$, $s=g-1$.
The reductive group $S=\operatorname{Sp}(W)\times\mathbb G_m$
acts by $t,t^{-1}$ on $a,b$. Both displayed conditions are
$S$-stable linear conditions, so its invariant projector
$e_S$ preserves them. Moreover
$e_{\operatorname{Sp}}e_S=e_{\operatorname{Sp}}$.
We calculate on the space of $S$-invariant line-tangent
derivations.

\paragraph{The complete invariant input.}
Regard the tensor $T_D$ as a word of length eight.
Let $\rho$ rotate its first letter to the last position,
let $S_8=1+\rho+\cdots+\rho^7$, and let
\[
 \mathcal D_n(v_1\cdots v_n)
       =[\cdots[[v_1,v_2],v_3],\ldots,v_n]
\]
be the left-normed Dynkin map. The conditions that $T_D$
be a special derivation are exactly
\[
 (1-\rho)T_D=0,\qquad
 (1\otimes\mathcal D_7)T_D=7T_D.
\]
Indeed the first equality is the kernel of the bracket
$V\otimes\mathbb L_7(V)\to\mathbb L_8(V)$, and the
second says that the last seven letters lie in the free
Lie algebra. The identity $\mathcal D_n^2=n\mathcal D_n$
gives the latter equivalence.

An $S$-invariant word contains the same number $m$ of
$a$'s and $b$'s, where $0\leq m\leq4$.
For $m>0$, its $b$-initial part is
$-b\otimes D(a)=-b\otimes[A,a]$ for a unique
$A\in\mathbb L_6(V)$ with $m-1$ copies each of $a,b$.
Uniqueness follows because the centralizer of $a$ in
degree six is zero. Summing the $b$-initial part over the
eight rotations counts every word $m$ times. Hence
\[
 T_D=\Psi_m(A):=-\frac1m S_8(b\otimes[A,a]).
\]
Conversely this formula gives the required derivation
exactly when
\begin{equation}\label{eq:line-input-constraints}
 (1\otimes\mathcal D_7)\Psi_m(A)=7\Psi_m(A),\qquad
 \operatorname{pr}_b\Psi_m(A)=-b\otimes[A,a].
\end{equation}
For $m=0$, every cyclic invariant tensor satisfying the
Dynkin condition has $D(a)=0$ and is allowed.

Here is a finite construction of every input in these
spaces. Place the indicated copies of $a,b$ in the six
slots for $A$, and match the remaining slots in all ways.
Each ordered pair represents
$\omega_W=\sum_{i=2}^g(x_i\otimes y_i-y_i\otimes x_i)$.
Apply $\mathcal D_6$, choose the independent columns of
this integer matrix, and impose
\eqref{eq:line-input-constraints} on their coefficients.
For $m=0$, start instead with all matchings of eight slots,
apply $S_8$, choose independent columns, and impose the
Dynkin condition. The symplectic invariant tensor theorem
\cite[Proposition~2.1, author p.~8]{MSS2015}
makes these lists spanning. They are linearly independent
as raw decorated diagrams in the range in use, since
$s\geq9$ and there are at most four pairs.
The exact row reductions have the following sizes:
\[
\begin{array}{c|rrrr}
m&\text{raw diagrams}&\text{parameter rank}&
       \text{constraint rank}&\text{kernel dimension}\\\hline
1&15&3&0&3\\
2&90&15&12&3\\
3&90&16&15&1\\
4&20&3&3&0\\
0&105&17&12&5
\end{array}
\]
Thus the full invariant line-tangent space has dimension
twelve. All entries of the reduction are given by the
two displayed word operations; no geometric relations
are imposed in this calculation.

\paragraph{The seven necessary scalar equations.}
Take two further symplectic pairs $(c,d),(e,f)$ in $W$.
For three vectors use the degree-one derivation
\[
 \delta_{uvw}(z)=\langle z,u\rangle[v,w]
          +\langle z,v\rangle[w,u]
          +\langle z,w\rangle[u,v].
\]
Its tensor is $-\operatorname{Alt}(u,v,w)$.
For a degree-seven special derivation $E$, write
$r_t(E)=\langle R(T_E),t\rangle$, with $R$ as in
Proposition~\ref{prop:johnson-seven-detector}.
The Enomoto--Satoh map first contracts the first two
slots of $T_E$ and takes the remaining seven slots
modulo cyclic rotation. It annihilates $J_{g,7}$
\cite[diagram (2), p.~640; \S7.1, p.~656]{EnomotoSatoh2014}.
For a matching $C$ of positions $1,\ldots,6$, define
on seven-letter cyclic words the scalar functional
\[
 \lambda_{C,t}(v_1\cdots v_7)
  =\sum_{j=0}^6
       \left(\prod_{(i,k)\in C}
          \langle v_{i+j},v_{k+j}\rangle\right)
                 \langle v_{7+j},t\rangle,
\]
where subscripts are read modulo seven. This is
well-defined on the cyclic quotient. Let $c_{C,t}(E)$
be its value after the Enomoto--Satoh contraction.
Put $C_1=(12)(35)(46)$ and $C_2=(12)(36)(45)$.
Define the seven-component column $M_7(D)$, in order, by
\[
\begin{pmatrix}
r_d([D,\delta_{abc}])\\
r_b([D,\delta_{acd}])\\
r_f([D,\delta_{cde}])\\
c_{C_1,d}([D,\delta_{abc}])\\
c_{C_2,d}([D,\delta_{abc}])\\
c_{C_1,b}([D,\delta_{acd}])\\
c_{C_1,f}([D,\delta_{cde}])
\end{pmatrix}.
\]
Every $\delta$ here belongs to $J_{g,1}=\bigwedge^3V$.
The second hypothesis of the lemma and the two detector
vanishings therefore give $M_7(e_SD)=0$.

\paragraph{The invariant quotient in degree six.}
Let $K(D)$ be the five scalar contractions of $T_D$
with the following ordered matchings:
\[
\begin{gathered}
(12)(34)(56)(78),\quad(12)(34)(57)(68),\quad
(12)(34)(58)(67),\\
(12)(36)(47)(58),\quad(12)(38)(46)(57).
\end{gathered}
\]
Morita--Sakasai--Suzuki's invariant calculation
\cite[Proposition~7.2, author pp.~26--27]{MSS2015}
says that $K$ restricts to an isomorphism
$\mathfrak h_{g,1}(6)^{\operatorname{Sp}}\to\mathbb Q^5$.
Their contraction convention differs by
$\operatorname{diag}(1,-1,1,-1,-1)$ from this ordered
pairing convention. Their two Johnson vectors have
coordinates $(0,0,0,-2,2)$ and $(4,-1,-4,4,-4)$
in their convention
\cite[proof of Theorem~7.5(i), author pp.~28--29]{MSS2015}.
Thus in the present convention their span is the kernel of
\[
 L=\begin{pmatrix}-1&4&0&0&0\\1&0&1&0&0\\0&0&0&1&1\end{pmatrix}.
\]
This identifies $K(J_{g,6}^{\operatorname{Sp}})=\ker L$.
Since each contraction is symplectically invariant,
$K(D)=K(e_{\operatorname{Sp}}D)$.

\paragraph{The polynomial identity.}
On all twelve inputs above, the literal contraction maps
satisfy
\begin{equation}\label{eq:geometric-polynomial-identity}
 N_0(s)M_7(s)=d_0(s)LK(s),\qquad
 d_0(s)=9(2s+3)(4s^2+6s-1).
\end{equation}
Here the entries of the $3$ by $7$ matrix $N_0$ are as follows:
\begin{align*}
(N_0)_{11}&=-2s(2s+3)(128s^4+528s^3+776s^2+465s+101),\\
(N_0)_{12}&=-4s(2s+3)(16s^4-48s^3-305s^2-402s-152),\\
(N_0)_{13}&=2s(s-1)(2s+3)(128s^3+464s^2+550s+199),\\
(N_0)_{14}&=-3s(576s^4+2720s^3+4684s^2+3546s+1029),\\
(N_0)_{15}&=-6s(s+2)(4s+3)(4s+5)(8s+7),\\
(N_0)_{16}&=3s(2s+1)(2s+3)(80s^2+240s+187),\\
(N_0)_{17}&=24s(s-1)(2s+3)(4s^2+6s-1).
\end{align*}
\begin{align*}
(N_0)_{21}&=2s(2s+1)^2(2s+3)(8s^2+25s+20),\\
(N_0)_{22}&=2s(2s+1)(2s+3)(4s^3-14s^2-70s-61),\\
(N_0)_{23}&=-4s(s-1)(2s+1)(2s+3)(8s^2+25s+20),\\
(N_0)_{24}&=3s(2s+3)(72s^3+232s^2+224s+69),\\
(N_0)_{25}&=6s(s+2)(2s+1)(2s+3)(8s+7),\\
(N_0)_{26}&=-3s(2s+3)(40s^3+140s^2+146s+37),\\
(N_0)_{27}&=-6s(s-1)(2s+3)(4s^2+6s-1).
\end{align*}
\begin{align*}
(N_0)_{31}&=-6s(2s+3)(16s^3+66s^2+80s+27),\\
(N_0)_{32}&=-6s(2s+3)(4s^3-12s^2-82s-81),\\
(N_0)_{33}&=6s(s-1)(2s+3)(16s^2+60s+53),\\
(N_0)_{34}&=-18s(s+2)(36s^2+102s+67),\\
(N_0)_{35}&=-36s(s+2)^2(8s+7),\\
(N_0)_{36}&=18s(s+2)(2s+3)(10s+13),\\
(N_0)_{37}&=0.
\end{align*}

We specify how the polynomial equality is evaluated.
Expand every Lie bracket by $uv-vu$, and every composition
of derivations by substitution in each letter. In the
symplectic tensor identification a substitution contracts
the output's first slot with the substituted slot. The
formula $T_{\delta_{uvw}}=-\operatorname{Alt}(u,v,w)$
fixes its sign. Thus all columns of $M_7$ and $K$ are
finite sums of fully paired decorated diagrams.
For each of the at most four copies of $\omega_W$, choose
one of its two orientations, recording its sign.
Each scalar pairing is zero unless the two letters have
opposite types $x,y$ and the same handle index; a pairing
in the order $y,x$ contributes a further minus sign.
Propagate these handle equalities. A $W$-index cannot
equal the handle of $a,b$, and the fixed handles of
$(a,b),(c,d),(e,f)$ are distinct. Each remaining free
handle has exactly $s$ choices, independently of all
other free handles. Its contribution is therefore $s$.
Summing the signed terms gives a polynomial in $s$,
including the signs and factors from closed contraction
cycles. Distinct cycles are allowed to choose the same
handle; no injectivity condition is inserted.

The row reductions in the table and this orientation
sum compute every entry of the two sides of
\eqref{eq:geometric-polynomial-identity} by rational
arithmetic. Subtraction gives three zero polynomials
for each of the twelve columns. The five pure-pair
columns of $K$ have determinant, up to the orientation
of the chosen kernel basis,
\[
 3538944s^5(s-2)(s-1)^4(s+1)^3
             (2s-1)(2s+1)^5(2s+3).
\]
This also checks their independence in the range used.
For a reproducible choice of columns, order the positions
of the $a$'s and $b$'s lexicographically and enumerate
matchings by pairing the first unused slot with each
later slot in increasing order. Use pivot columns in
reduced row echelon form for each parameter basis and
the associated free-coordinate basis for each kernel;
clear denominators and make the first nonzero coefficient
positive. The program \texttt{verify\_single\_line.py},
cited through the \href{\workbenchurl}{workbench},
performs exactly these word operations and orientation
sums, and gives all thirty-six zero differences.

Since $s\geq9$, $d_0(s)$ is nonzero. For example
$d_0(3+t)=9(2t+9)(4t^2+30t+53)>0$ for $t\geq0$.
The equality $M_7(e_SD)=0$ and
\eqref{eq:geometric-polynomial-identity} imply
$LK(e_SD)=0$. Hence
$K(e_{\operatorname{Sp}}D)\in K(J_{g,6}^{\operatorname{Sp}})$.
Injectivity of $K$ on symplectic invariants proves
$e_{\operatorname{Sp}}D\in J_{g,6}$.
Every symplectically trivial quotient has
$q(D)=q(e_{\operatorname{Sp}}D)$, giving the last assertion.
\end{proof}
